% Generated by code/interviews-1/03-speed-drills/python/gen_drill.py. Do not edit. French copy by draft/pilot-quant-fr/ivgen_fr.py.
\ivfamily{Trois chiffres fois deux chiffres}{2, 13, 14, 15, 20, 24, 27, 33}{Ne pas finir le calcul. Vérifier le chiffre des unités, puis l'ordre de grandeur, puis la somme des chiffres. Pour $378 \times 46$, le chiffre des unités est celui de $8 \times 6$, et l'ordre de grandeur est voisin de $400 \times 45$.}
\ivfamily{Carrés et quasi-carrés}{16, 17, 19, 21, 26}{$67^2$ finit par 9, et est proche de $70 \times 64$. $104 \times 96$ vaut $10000 - 16$.}
\ivfamily{Division exacte}{1, 7, 18, 25, 31, 36}{Vérifier en calcul mental~: le quotient multiplié par le diviseur doit finir par le bon chiffre. Son ordre de grandeur vient de l'arrondi du diviseur.}
\ivfamily{Pourcentages}{6, 10, 12, 23, 28, 30}{$17\%$ de 640 vaut $10\% + 5\% + 2\%$, soit $64 + 32 + 12.8$.}
\ivfamily{Produits de décimaux}{3, 4, 8, 11, 29, 35}{Compter d'abord les décimales. Les propositions ne diffèrent alors que par des chiffres qu'on peut tester~: $3.4 \times 1.7$ a deux décimales et finit par 8.}
\ivfamily{Fractions d'un nombre}{5, 9, 22, 32, 34}{Diviser d'abord par le dénominateur, puis multiplier par le numérateur. Pour $\tfrac78$, retrancher un huitième.}
